\bookchapter{第六十一章：Swiftにすばやく入門——Apple開発の第一言語}

\chapterintro{この章から、いよいよ言語を学びます。緊張しないでください——Swiftは「初心者にやさしい」で有名ですし、文法書を丸暗記する必要もありません。文法の20％をまず押さえれば、日常のコードの80％は読めて書けます。残りはプロジェクトを進めながら調べればOK。私たちの学習の原則を忘れずに：理解する＞暗記する、使える＞完璧。}

\booksection{変数と型：データに名前をつける}

プログラムとは「データ＋ロジック」です。変数は、データに名前をつけてしまっておくもの。Swiftではletとvarで宣言します：

\begin{lstlisting}[language=Swift]
let name = "たろう"        // let：定数（デフォルトはこれがおすすめ）
var age = 18             // var：変数（あとから変えられる）

age = 19                 // <i class="fas fa-check-circle mr-1"></i> varは変えられる
// name = "はなこ"         // <i class="fas fa-times-circle mr-1"></i> letは変えられない。コンパイルエラー

let score: Int = 100     // 型を明示的に書いてもいい
let pi = 3.14            // 書かなくても自動で推論される

print("こんにちは、\(name)")   // 文字列補間：バックスラッシュ + 括弧
\end{lstlisting}

よく使う型は、この1枚の表で覚えましょう：

\begin{center}
\begin{tabularx}{\linewidth}{|X|X|X|}
\hline
\textbf{型} & \textbf{何を入れるか} & \textbf{例} \tabularnewline
\hline
Int & 整数 & 42 \tabularnewline
Double & 小数 & 3.14 \tabularnewline
String & 文字列 & “こんにちは” \tabularnewline
Bool & 真／偽 & true / false \tabularnewline
{}[Int]（配列） & データのまとまり & {}[1, 2, 3] \tabularnewline
\hline
\end{tabularx}
\end{center}

\begin{center}\small\color{gray}Swiftのよく使う基本型\end{center}

\begin{tcbinfo}{ワンポイント：let優先の原則}
Swiftコミュニティの鉄則：デフォルトはlet、本当に必要なときだけvar。letで足りる場所にvarを使うと、Xcodeが黄色い文字で警告してくれます。この習慣があるだけで、「変数が意図せず書き換わる」系のバグの一大カテゴリを減らせます——変化の少ないコードほど、頭の中で追いやすいのです。

\end{tcbinfo}

\booksection{オプショナル（Optional）：Swiftの「クラッシュ防止のお守り」}

これはSwiftでいちばん重要な概念で、ほかにはありません。ひとことで言えば、値が「存在しない」かもしれないとき、Swiftはあなたに処理させる——「ヌルポインタクラッシュ」という業界の悪夢を、言語レベルで消し去ったのです。

\begin{lstlisting}[language=Swift]
var name: String = "たろう"   // 普通の型：nilは絶対に許されない
// name = nil                // <i class="fas fa-times-circle mr-1"></i> コンパイラが直接止めてくれる

var nickname: String? = nil  // 疑問符を付けると：nilを許す（オプショナル）
nickname = "たーくん"

// オプショナルを使うときは「ほどく」必要がある：
if let n = nickname {        // if let：安全にほどく
    print("こんにちは、\(n)")
} else {
    print("まだニックネームはない")
}

let len = nickname?.count ?? 0   // ?? でデフォルト値（nil結合演算子）
print(len)
\end{lstlisting}

\begin{tcbwarn}{注意：!は「危険な操作」}
コードにnickname!（強制アンラップ）と書くのは、コンパイラに「nilじゃないって保証する！」と叫ぶのと同じ——外れたらアプリはその場でクラッシュします。初心者のうちはif let / ??を使い、!は使わないのが原則。Xcodeがコードを補完するとき、たまに!を勝手に付けてくるので、直すのを忘れずに。

\end{tcbwarn}

\booksection{関数とクロージャ：ロジックを「ボタン」にまとめる}

関数は名前のついたロジックのかたまりで、リモコンのボタンのようなもの。クロージャ（Closure）は「名前のない関数」で、その場でほかに渡せます：

\begin{lstlisting}[language=Swift]
// 関数：名前が先、引数と戻り値の型が後
func greet(name: String) -> String {
    return "こんにちは、\(name)！"
}
print(greet(name: "たろう"))     // 呼び出すときは引数名を付ける

// クロージャ：その場で定義。手軽でストレート
let numbers = [1, 2, 3, 4]
let doubled = numbers.map { $0 * 2 }   // $0 = 現在の要素
// [2, 4, 6, 8]

let big = numbers.filter { $0 > 2 }    // 絞り込み
// [3, 4]
\end{lstlisting}

SwiftUIではクロージャが至るところに登場します。ボタンのButton(“押して”) \{ ... \}の後ろの中かっこもクロージャです。map、filter、sortedという「コレクションのコンボ」はしっかり覚えておきましょう——この先のリスト処理の90％は、この組み合わせです。

\booksection{構造体とプロトコル：Swiftの「レゴ」システム}

Swiftの公式のお好みははっきりしています。基本は構造体（struct）で、継承が必要なときだけクラス（class）。構造体は「値型」——コピーすればそれぞれ独立して動くので、ある種のバグが生まれにくいのです：

\begin{lstlisting}[language=Swift]
// struct：データを入れる箱
struct Record {
    var id = UUID()
    var amount: Double
    var category: String
    var date: Date

    // structは「メンバー逐一イニシャライザ」付き：すぐ作れる
}
let r = Record(amount: 25.5, category: "食費", date: Date())

// protocol：「どんな能力が必須か」の約束
protocol Identifiable {
    var id: UUID { get }
}
struct Record: Identifiable { ... }   // protocolへの準拠 = 要件の実装を約束すること
\end{lstlisting}

SwiftUIではほぼすべてがstructです（Viewはプロトコル）。いまはピンとこなくて大丈夫。第六十二章で毎日のように顔を合わせます——Swiftの世界では、まずstructを受け入れてしまいましょう。

\booksection{ゼロからの初心者のための文法チートシート}

\begin{center}
\begin{tabularx}{\linewidth}{|X|X|}
\hline
\textbf{やりたいこと} & \textbf{書き方} \tabularnewline
\hline
定数／変数の定義 & \texttt{let x = 1}/\texttt{var x = 1} \tabularnewline
文字列への値の埋め込み & “こんにちは、\textbackslash{}(name)” \tabularnewline
関数の定義 & func 名前(引数: 型) -> 戻り値の型 \{ ... \} \tabularnewline
条件分岐 & if 条件 \{ ... \} else \{ ... \} \tabularnewline
繰り返し & for x in 配列 \{ ... \} \tabularnewline
安全なアンラップ & if let x = オプショナル値 \{ ... \} \tabularnewline
デフォルト値 & x ?? デフォルト値 \tabularnewline
構造体の定義 & \texttt{struct X \{ var a: Int; var b: String \}} \tabularnewline
クロージャ & \texttt{\{ \$0 * 2 \}} \tabularnewline
\hline
\end{tabularx}
\end{center}

\begin{center}\small\color{gray}Swift文法チートシート（初心者版）\end{center}

\begin{tcbpractice}{やってみよう：PlaygroundsでSwiftを15分遊ぶ}
MacがあればXcode Playgroundを開きましょう。なければiPadのSwift Playgroundsや、ブラウザ版のプレイグラウンドでOK。この章の各コードをすべて実行し、数字や文字をいじって、わざとエラーをいくつか起こしてXcodeがどう教えてくれるか見てみてください。この「遊んでから学ぶ」15分のおかげで、第六十二章に入るときのハードルは半分になります。

\end{tcbpractice}

\begin{tcbsummary}{本章のまとめ}
\begin{itemize}
  \item letを優先、varは慎重に。型は自動推論、コロンのあとに明示宣言してもOK
  \item オプショナルはSwiftの魂：空かもしれないなら必ず処理する。if let / ??はお守り、!は危険な操作
  \item 関数呼び出しには引数名をつける。クロージャ＝名前のない関数で、SwiftUIでは至るところに登場
  \item structを優先し、プロトコル（Protocol）に従う。SwiftUIではほぼすべてがstruct
  \item 文法の20％で仕事の80％：まず使えるようになって、それから極める
\end{itemize}

\end{tcbsummary}
